\chapter{Complexity — Primality Math}
%%% generated by the blueprint pipeline from TCSlib.Complexity.Randomized.PrimalityMath
\section{Overview}

Declarations of the Lean module \texttt{TCSlib.Complexity.Randomized.PrimalityMath}.

\section{Declarations}

\begin{definition}[Set of prime numbers]
\lean{Randomized.PRIMES}
\statementsource{Vad12}{
  pdf-bf2ce5422088-p020-b001,
  pdf-bf2ce5422088-p046-b002
}
\statementsource{arora-barak-ch07-randomized-computation}{pdf-26d95080dbde-p004-b001}
\label{Randomized.PRIMES}
\leanok
The set of prime numbers is $\{N\in\mathbb{N}:N\text{ is prime}\}$.
\end{definition}

\begin{definition}[Quadratic-residuosity symbol]
\lean{Randomized.qrSymbol}
\statementsource{arora-barak-ch07-randomized-computation}{pdf-26d95080dbde-p004-b002}
\label{Randomized.qrSymbol}
\leanok
For natural numbers $N$ and $A$, define the quadratic-residuosity symbol to be $0$ when
$A$ and $N$ are not coprime; when they are coprime, it is $1$ if the residue class of
$A$ in $\mathbb{Z}/N\mathbb{Z}$ is a square, and $-1$ otherwise.
\end{definition}

\begin{theorem}[Three Possible Values of the Quadratic Residue Symbol]
\lean{Randomized.qrSymbol_trichotomy}
\label{Randomized.qrSymbol_trichotomy}
\leanok
\uses{Randomized.qrSymbol}
\difficulty{3}
For any natural numbers $N$ and $A$, define the quadratic residue symbol modulo $N$ to
be $0$ when $A$ is not coprime to $N$, $1$ when the residue of $A$ modulo $N$ is a
square, and $-1$ otherwise. This symbol is one of $0$, $1$, or $-1$.
\end{theorem}

\begin{theorem}[Coprimality and nonzero residues]
\lean{Randomized.coprime_iff_natCast_ne_zero}
\label{Randomized.coprime_iff_natCast_ne_zero}
\leanok
\difficulty{3}
If $p$ is prime and $A$ is a natural number, then $A$ is coprime to $p$ if and only if
its residue modulo $p$ is nonzero.
\end{theorem}

\begin{theorem}[Qr symbol eq legendre sym]
\lean{Randomized.qrSymbol_eq_legendreSym}
\statementsource{arora-barak-ch07-randomized-computation}{pdf-26d95080dbde-p004-b002}
\label{Randomized.qrSymbol_eq_legendreSym}
\leanok
\uses{Randomized.coprime_iff_natCast_ne_zero, Randomized.qrSymbol}
\difficulty{5}
% dataset-pipeline: stub -- prose pending build_informal_statements.py
Statement pending regeneration.
\end{theorem}

\begin{theorem}[Jacobi sym eq prod qr symbol]
\lean{Randomized.jacobiSym_eq_prod_qrSymbol}
\statementsource{arora-barak-ch07-randomized-computation}{pdf-26d95080dbde-p004-b002}
\label{Randomized.jacobiSym_eq_prod_qrSymbol}
\leanok
\uses{Randomized.qrSymbol, Randomized.qrSymbol_eq_legendreSym}
\difficulty{4}
% dataset-pipeline: stub -- prose pending build_informal_statements.py
Statement pending regeneration.
\end{theorem}

\begin{theorem}[Euler’s Criterion for the Quadratic Residue Symbol]
\lean{Randomized.qrSymbol_eq_pow}
\proofsource{arora-barak-ch07-randomized-computation}{pdf-26d95080dbde-p004-b002}
\label{Randomized.qrSymbol_eq_pow}
\leanok
\uses{Randomized.qrSymbol, Randomized.qrSymbol_eq_legendreSym}
\difficulty{3}
\proofstep
  {Randomized.qrSymbol}
  {direct}
  {arora-barak-ch07-randomized-computation}
  {pdf-26d95080dbde-p004-b002}
\proofstep
  {Randomized.coprime_iff_natCast_ne_zero}
  {context}
  {arora-barak-ch07-randomized-computation}
  {pdf-26d95080dbde-p004-b002}
\proofstep
  {Randomized.qrSymbol_eq_legendreSym}
  {direct}
  {arora-barak-ch07-randomized-computation}
  {pdf-26d95080dbde-p004-b002}
\proofstep
  {Randomized.IsEulerLiar}
  {context}
  {arora-barak-ch07-randomized-computation}
  {pdf-26d95080dbde-p004-b002}
\proofstep
  {Randomized.instDecidableIsEulerLiar}
  {context}
  {arora-barak-ch07-randomized-computation}
  {pdf-26d95080dbde-p004-b002}
Let $p$ be an odd prime and let $A$ be a natural number. Define the quadratic residue
symbol modulo $p$ to be $0$ when $A$ is not coprime to $p$, $1$ when the residue of $A$
modulo $p$ is a square, and $-1$ otherwise. Its value modulo $p$ equals $A^{(p-1)/2}$
modulo $p$.
\end{theorem}

\begin{definition}[Euler liars]
\lean{Randomized.IsEulerLiar}
\statementsource{arora-barak-ch07-randomized-computation}{pdf-26d95080dbde-p004-b002}
\label{Randomized.IsEulerLiar}
\leanok
For natural numbers $N$ and $A$, say that $A$ is an Euler liar for $N$ if $A$ and $N$
are coprime and, in $\mathbb{Z}/N\mathbb{Z}$, the residue class of the Jacobi symbol
$(A/N)$ equals the $(N-1)/2$-th power of the residue class of $A$.
\end{definition}

\begin{theorem}[Remainder characterization of Euler liars]
\lean{Randomized.isEulerLiar_iff_emod}
\label{Randomized.isEulerLiar_iff_emod}
\leanok
\uses{Randomized.IsEulerLiar}
\difficulty{3}
For natural numbers $N$ and $A$, the numbers $A$ is an Euler liar for $N$ if and only if
$\gcd(A,N)=1$ and the least nonnegative remainder of the Jacobi symbol $(A/N)$ upon
division by $N$ equals the remainder of $A^{(N-1)/2}$ upon division by $N$.
\end{theorem}

\begin{theorem}[Primes are Euler liars]
\lean{Randomized.isEulerLiar_of_prime}
\proofsource{arora-barak-ch07-randomized-computation}{pdf-26d95080dbde-p004-b002}
\label{Randomized.isEulerLiar_of_prime}
\leanok
\uses{Randomized.IsEulerLiar, Randomized.qrSymbol_eq_legendreSym, Randomized.qrSymbol_eq_pow}
\difficulty{3}
\proofstep
  {Randomized.qrSymbol}
  {direct}
  {arora-barak-ch07-randomized-computation}
  {pdf-26d95080dbde-p004-b002}
\proofstep
  {Randomized.coprime_iff_natCast_ne_zero}
  {context}
  {arora-barak-ch07-randomized-computation}
  {pdf-26d95080dbde-p004-b002}
\proofstep
  {Randomized.qrSymbol_eq_legendreSym}
  {context}
  {arora-barak-ch07-randomized-computation}
  {pdf-26d95080dbde-p004-b002}
\proofstep
  {Randomized.qrSymbol_eq_pow}
  {direct}
  {arora-barak-ch07-randomized-computation}
  {pdf-26d95080dbde-p004-b002}
\proofstep
  {Randomized.IsEulerLiar}
  {direct}
  {arora-barak-ch07-randomized-computation}
  {pdf-26d95080dbde-p004-b002}
\proofstep
  {Randomized.instDecidableIsEulerLiar}
  {context}
  {arora-barak-ch07-randomized-computation}
  {pdf-26d95080dbde-p004-b002}
Let $p$ be a prime different from $2$, and let $A$ be a natural number coprime to $p$.
Then $A$ is an Euler liar for $p$: the Jacobi symbol $(A/p)$, viewed modulo $p$, equals
$A^{(p-1)/2}$ modulo $p$.
\end{theorem}

\begin{theorem}[Euler liar for a bounded prime residue]
\lean{Randomized.isEulerLiar_of_prime_of_mem_Ico}
\label{Randomized.isEulerLiar_of_prime_of_mem_Ico}
\leanok
\uses{Randomized.IsEulerLiar, Randomized.coprime_iff_natCast_ne_zero, Randomized.isEulerLiar_of_prime}
\difficulty{3}
Let $p$ be a prime different from $2$, and let $A$ be a natural number with $1\le A<p$.
Then $A$ is an Euler liar for $p$: the Jacobi symbol $(A/p)$, viewed modulo $p$, equals
$A^{(p-1)/2}$ modulo $p$.
\end{theorem}

\begin{definition}[Jacobi symbol in a residue ring]
\lean{Randomized.jacobiSymZMod}
\label{Randomized.jacobiSymZMod}
\leanok
For each natural number $N$ and each residue class $x$ modulo $N$, let $r$ be the least
nonnegative integer representative of $x$. The defined residue class modulo $N$ is the
Jacobi symbol of $r$ with modulus $N$, reduced modulo $N$.
\end{definition}

\begin{theorem}[Jacobi symbol on natural casts]
\lean{Randomized.jacobiSymZMod_natCast}
\label{Randomized.jacobiSymZMod_natCast}
\leanok
\uses{Randomized.jacobiSymZMod}
\difficulty{3}
For every positive natural number $N$ and every natural number $A$, the Jacobi symbol
viewed as a residue class modulo $N$ at the residue class of $A$ equals the residue
class modulo $N$ of the integer-valued Jacobi symbol of $A$ and $N$.
\end{theorem}

\begin{theorem}[Multiplicativity of the Jacobi symbol]
\lean{Randomized.jacobiSymZMod_mul}
\label{Randomized.jacobiSymZMod_mul}
\leanok
\uses{Randomized.jacobiSymZMod}
\difficulty{3}
For every positive natural number $N$ and all residue classes $x,y$ modulo $N$, the
Jacobi symbol, viewed as a residue class modulo $N$, is multiplicative: its value at
$xy$ equals the product of its values at $x$ and $y$.
\end{theorem}

\begin{theorem}[Jacobi symbol at one]
\lean{Randomized.jacobiSymZMod_one}
\label{Randomized.jacobiSymZMod_one}
\leanok
\uses{Randomized.jacobiSymZMod, Randomized.jacobiSymZMod_natCast}
\difficulty{3}
For every positive natural number $N$, the Jacobi symbol modulo $N$, viewed as a residue
class in $\mathbb{Z}/N\mathbb{Z}$, takes the value $1$ at the multiplicative identity.
\end{theorem}

\begin{definition}[Units satisfying the Euler–Jacobi relation]
\lean{Randomized.eulerLiarSubgroup}
\label{Randomized.eulerLiarSubgroup}
\leanok
\uses{Randomized.jacobiSymZMod, Randomized.jacobiSymZMod_mul, Randomized.jacobiSymZMod_one}
For a natural number $N$ with $N\ne 0$, consider the units modulo $N$. The defined
subgroup consists of those units $u$ for which the Jacobi symbol modulo $N$ of the
residue represented by $u$ equals $u^{(N-1)/2}$ in the residue ring modulo $N$.
\end{definition}

\begin{theorem}[Membership in the Euler Liar Subgroup]
\lean{Randomized.mem_eulerLiarSubgroup}
\label{Randomized.mem_eulerLiarSubgroup}
\leanok
\uses{Randomized.eulerLiarSubgroup, Randomized.jacobiSymZMod}
\difficulty{1}
Let $N$ be a nonzero natural number and let $u$ be a unit modulo $N$. Then $u$ belongs
to the Euler liar subgroup precisely when the Jacobi symbol of its residue modulo $N$,
viewed in $\mathbb{Z}/N\mathbb{Z}$, equals $u^{(N-1)/2}$ in $\mathbb{Z}/N\mathbb{Z}$.
\end{theorem}

\begin{theorem}[Exponentiation of an Euler-Liar Unit]
\lean{Randomized.pow_sub_one_eq_one_of_mem_eulerLiarSubgroup}
\label{Randomized.pow_sub_one_eq_one_of_mem_eulerLiarSubgroup}
\leanok
\uses{Randomized.eulerLiarSubgroup, Randomized.jacobiSymZMod, Randomized.mem_eulerLiarSubgroup}
\difficulty{4}
Let $N>1$ be an odd natural number, and let $u$ be a unit modulo $N$ whose residue class
satisfies $\left(\frac{u}{N}\right)=u^{(N-1)/2}$ in $\mathbb{Z}/N\mathbb{Z}$, where the
left side is the Jacobi symbol interpreted modulo $N$. Then $u^{N-1}=1$ in
$\mathbb{Z}/N\mathbb{Z}$.
\end{theorem}

\begin{theorem}[A Unit Outside the Euler-Liar Subgroup]
\lean{Randomized.exists_not_mem_of_sq_dvd}
\label{Randomized.exists_not_mem_of_sq_dvd}
\leanok
\uses{Randomized.eulerLiarSubgroup, Randomized.pow_sub_one_eq_one_of_mem_eulerLiarSubgroup}
\difficulty{5}
Let $N>1$ be odd, and let $p$ be a prime such that $p^2$ divides $N$. There is a unit
$u$ modulo $N$ outside the subgroup of units satisfying
$\left(\frac{u}{N}\right)=u^{(N-1)/2}$ in $\mathbb{Z}/N\mathbb{Z}$, where the left side
is the Jacobi symbol, viewed modulo $N$.
\end{theorem}

\begin{theorem}[A Unit Outside the Euler-Liar Subgroup]
\lean{Randomized.exists_not_mem_of_not_sq_dvd}
\label{Randomized.exists_not_mem_of_not_sq_dvd}
\leanok
\uses{Randomized.coprime_iff_natCast_ne_zero, Randomized.eulerLiarSubgroup, Randomized.jacobiSymZMod_natCast, Randomized.mem_eulerLiarSubgroup}
\difficulty{6}
Let $N>1$ be an odd nonprime integer, and suppose no prime square divides $N$. There is
a unit $u$ modulo $N$ outside the subgroup of units satisfying
$\left(\frac{u}{N}\right)=u^{(N-1)/2}$ in $\mathbb{Z}/N\mathbb{Z}$, where the left side
is the Jacobi symbol, viewed modulo $N$.
\end{theorem}

\begin{theorem}[The Euler-liar subgroup is proper]
\lean{Randomized.eulerLiarSubgroup_ne_top}
\label{Randomized.eulerLiarSubgroup_ne_top}
\leanok
\uses{Randomized.eulerLiarSubgroup, Randomized.exists_not_mem_of_not_sq_dvd, Randomized.exists_not_mem_of_sq_dvd}
\difficulty{4}
Let $N$ be a nonzero natural number with $1<N$, assume that $N$ is odd and not prime.
Then the subgroup of units modulo $N$ whose Jacobi symbol equals their $(N-1)/2$-th
power is a proper subgroup of the units modulo $N$.
\end{theorem}

\begin{theorem}[Two mul card euler liar subgroup le]
\lean{Randomized.two_mul_card_eulerLiarSubgroup_le}
\statementsource{arora-barak-ch07-randomized-computation}{pdf-26d95080dbde-p004-b002}
\label{Randomized.two_mul_card_eulerLiarSubgroup_le}
\leanok
\uses{Randomized.eulerLiarSubgroup, Randomized.eulerLiarSubgroup_ne_top}
\difficulty{3}
% dataset-pipeline: stub -- prose pending build_informal_statements.py
Statement pending regeneration.
\end{theorem}

\begin{theorem}[Euler liars among residues]
\lean{Randomized.two_mul_card_filter_isEulerLiar_le}
\proofsource{arora-barak-ch07-randomized-computation}{pdf-26d95080dbde-p004-b002}
\label{Randomized.two_mul_card_filter_isEulerLiar_le}
\leanok
\uses{Randomized.IsEulerLiar, Randomized.eulerLiarSubgroup, Randomized.jacobiSymZMod_natCast, Randomized.mem_eulerLiarSubgroup, Randomized.two_mul_card_eulerLiarSubgroup_le}
\difficulty{5}
\proofstep
  {Randomized.coprime_iff_natCast_ne_zero}
  {context}
  {arora-barak-ch07-randomized-computation}
  {pdf-26d95080dbde-p004-b002}
\proofstep
  {Randomized.IsEulerLiar}
  {direct}
  {arora-barak-ch07-randomized-computation}
  {pdf-26d95080dbde-p004-b002}
\proofstep
  {Randomized.instDecidableIsEulerLiar}
  {context}
  {arora-barak-ch07-randomized-computation}
  {pdf-26d95080dbde-p004-b002}
\proofstep
  {Randomized.jacobiSymZMod}
  {context}
  {arora-barak-ch07-randomized-computation}
  {pdf-26d95080dbde-p004-b002}
\proofstep
  {Randomized.jacobiSymZMod_natCast}
  {context}
  {arora-barak-ch07-randomized-computation}
  {pdf-26d95080dbde-p004-b002}
\proofstep
  {Randomized.jacobiSymZMod_mul}
  {context}
  {arora-barak-ch07-randomized-computation}
  {pdf-26d95080dbde-p004-b002}
\proofstep
  {Randomized.jacobiSymZMod_one}
  {context}
  {arora-barak-ch07-randomized-computation}
  {pdf-26d95080dbde-p004-b002}
\proofstep
  {Randomized.eulerLiarSubgroup}
  {context}
  {arora-barak-ch07-randomized-computation}
  {pdf-26d95080dbde-p004-b002}
\proofstep
  {Randomized.mem_eulerLiarSubgroup}
  {context}
  {arora-barak-ch07-randomized-computation}
  {pdf-26d95080dbde-p004-b002}
\proofstep
  {Randomized.pow_sub_one_eq_one_of_mem_eulerLiarSubgroup}
  {context}
  {arora-barak-ch07-randomized-computation}
  {pdf-26d95080dbde-p004-b002}
\proofstep
  {Randomized.exists_not_mem_of_sq_dvd}
  {context}
  {arora-barak-ch07-randomized-computation}
  {pdf-26d95080dbde-p004-b002}
\proofstep
  {Randomized.exists_not_mem_of_not_sq_dvd}
  {context}
  {arora-barak-ch07-randomized-computation}
  {pdf-26d95080dbde-p004-b002}
\proofstep
  {Randomized.eulerLiarSubgroup_ne_top}
  {context}
  {arora-barak-ch07-randomized-computation}
  {pdf-26d95080dbde-p004-b002}
\proofstep
  {Randomized.two_mul_card_eulerLiarSubgroup_le}
  {direct}
  {arora-barak-ch07-randomized-computation}
  {pdf-26d95080dbde-p004-b002}
Let $N>1$ be an odd composite integer. Among the integers $A$ with $1\le A<N$ that are
coprime to $N$ and satisfy $\left(\frac{A}{N}\right)\equiv A^{(N-1)/2}\pmod N$, at most
$(N-1)/2$ have this property.
\end{theorem}

\begin{definition}[Euler test sample]
\lean{Randomized.eulerSample}
\label{Randomized.eulerSample}
\leanok
\uses{Complexity.LogProg.bitsVal}
For a natural number $N$ and a Boolean list $r$, the defined sample is one plus the
remainder when the integer encoded by $r$ in binary (with the first list entry as the
least significant bit) is divided by $N-1$.
\end{definition}

\begin{theorem}[Range of the Euler sample]
\lean{Randomized.eulerSample_mem_Ico}
\label{Randomized.eulerSample_mem_Ico}
\leanok
\uses{Complexity.LogProg.bitsVal, Randomized.eulerSample}
\difficulty{3}
Let $N$ be a natural number with $1<N$, and let $r$ be a finite binary string. The Euler
sample obtained from $N$ and $r$ lies in the interval of natural numbers $[1,N)$.
\end{theorem}

\begin{theorem}[Rand prob is euler liar le]
\lean{Randomized.randProb_isEulerLiar_le}
\statementsource{arora-barak-ch07-randomized-computation}{pdf-26d95080dbde-p004-b002}
\label{Randomized.randProb_isEulerLiar_le}
\leanok
\uses{Complexity.LogProg.bitsVal, Complexity.LogProg.bitsVal_lt_two_pow, Complexity.LogProg.eq_of_bitsVal_eq, Randomized.IsEulerLiar, Randomized.eulerSample, Randomized.eulerSample_mem_Ico, Randomized.randProb, Randomized.two_mul_card_filter_isEulerLiar_le}
\difficulty{6}
% dataset-pipeline: stub -- prose pending build_informal_statements.py
Statement pending regeneration.
\end{theorem}

\begin{definition}[Repeated Euler–Jacobi liar test]
\lean{Randomized.eulerLiarRounds}
\label{Randomized.eulerLiarRounds}
\leanok
\uses{Randomized.IsEulerLiar, Randomized.eulerSample}
For natural numbers $N$ and $\ell$, the test on $k$ rounds and a Boolean list $r$
returns true when $k=0$. For $k+1$ rounds, it returns true exactly when the sample
formed from the first $\ell$ entries of $r$ is an Euler liar for $N$ and the remaining
list passes the test for $k$ rounds.
\end{definition}

\begin{theorem}[Two rounds depend only on the first two blocks]
\lean{Randomized.eulerLiarRounds_two_take}
\label{Randomized.eulerLiarRounds_two_take}
\leanok
\uses{Randomized.eulerLiarRounds}
\difficulty{1}
For all natural numbers $N$ and $\ell$ and every binary string $r$, the result of two
Euler-liar rounds using blocks of length $\ell$ is unchanged when the random string is
truncated to its first $2\ell$ bits.
\end{theorem}

\begin{theorem}[Euler liar rounds of prime]
\lean{Randomized.eulerLiarRounds_of_prime}
\statementsource{arora-barak-ch07-randomized-computation}{pdf-26d95080dbde-p004-b002}
\label{Randomized.eulerLiarRounds_of_prime}
\leanok
\uses{Randomized.eulerLiarRounds, Randomized.eulerSample_mem_Ico, Randomized.isEulerLiar_of_prime_of_mem_Ico}
\difficulty{4}
% dataset-pipeline: stub -- prose pending build_informal_statements.py
Statement pending regeneration.
\end{theorem}

\begin{theorem}[Euler liar probability over repeated rounds]
\lean{Randomized.randProb_eulerLiarRounds_le}
\label{Randomized.randProb_eulerLiarRounds_le}
\leanok
\uses{Randomized.IsEulerLiar, Randomized.eulerLiarRounds, Randomized.eulerSample, Randomized.randProb, Randomized.randProb_congr, Randomized.randProb_isEulerLiar_le, Randomized.randProb_le_one, Randomized.randProb_nonneg, Randomized.randProb_split}
\difficulty{5}
Let $N>1$ be an odd composite natural number, and let $\ell,k$ be natural numbers. In
each of $k$ rounds, use a block of $\ell$ uniformly random bits to form a candidate by
interpreting the block as a binary integer, reducing it modulo $N-1$, and adding $1$;
the rounds accept when every candidate is an Euler liar to base $N$. The probability of
acceptance using $\ell k$ uniformly random bits is at most
$\left(\frac12+\frac{N-1}{2^{\ell+1}}\right)^k$.
\end{theorem}

\begin{theorem}[Two-round Euler liar probability bound]
\lean{Randomized.randProb_eulerLiarRounds_two_le}
\label{Randomized.randProb_eulerLiarRounds_two_le}
\leanok
\uses{Randomized.eulerLiarRounds, Randomized.randProb, Randomized.randProb_eulerLiarRounds_le}
\difficulty{4}
Let $N>1$ be an odd composite natural number, and let $\ell$ be a natural number such
that $8N\le 2^\ell$. Run two Euler liar tests for $N$, each using an independent block
of $\ell$ uniformly random bits to form its candidate by interpreting the block as a
binary integer, reducing it modulo $N-1$, and adding $1$. The probability that both
rounds accept, using $2\ell$ uniformly random bits, is at most $\frac13$.
\end{theorem}
